\chapter{Противоречия капиталистической экономики и их проявление в рамках кризиса}
\label{app:dialectica}

\subsection{Капиталистическая экономика как внутренне противоречивая система}
Поскольку капиталистическая экономика представляет собой исторически конкретную форму организации общественного воспроизводства, ее анализ требует выделения содержащихся в самой ее основе фундаментальных противоречий, раскрывающих сущность, динамику и исторические пределы ее развития, задающих тем самым внутреннюю логику ее функционирования и обуславливающих закономерно возникающие в ходе ее движения диспропорции, кризисы и качественные переломы, формирующие основу предпосылок ее дальнейшего отрицания.

В рамках капиталистической экономики фундаментальные противоречия всех ее составных элементов возможно классифицировать по двум основные категориям:
\begin{itemize}
    \item \textit{Главные противоречия}, непосредственно определяющие саму капиталистическую экономику, не способные быть разрешенными в рамках системы общественных отношений капиталистического строя и составляющие основу более конкретных второстепенных противоречий;
    \item \textit{Второстепенные противоречия}, лежащие на поверхности и представляющие собой более конкретную производную от главных противоречий форму, частично снимающиеся посредством обусловленного ими насильственного разрешения накопившихся количественных и структурных диспропорций в рамках кризиса капиталистической экономики и вновь воспроизводящиеся в дальнейшем на более высокой стадии ее развития;
\end{itemize}

Материальной базой капиталистической экономики является исторически развернутый капиталистический способ производства, включающий в себя следующие основополагающие элементы:
\begin{itemize}
    \item \textit{Товар}, представляющий собой исторически обусловленную овеществленную форму продукта труда, первоначально возникшую с разложением первобытного строя на этапе обмена излишков между различными общинами и племенами, получившую ограниченное распространение в докапиталистическую эпоху с развитием ремесленного производства и торговли в рамках крупных городов и достигшую всеобщего характера опосредующего производственные отношения центрального элемента капиталистического способа производства;
    \item \textit{Деньги}, представляющие собой всеобщий эквивалент стоимости, изначально возникший на базе стихийно выделившихся в процессе развития товарного обращения специфических товаров в качестве всеобщей формы стоимости и впоследствии превратившийся в связи с углублением капиталистических отношений и объективной ограниченностью денежных товаров в лишенный товарного содержания и сохраняющий в себе все исходные функции товарных денег номинальный знак стоимости; законодательно закрепленный и монополизированный государством в качестве законного платежного средства и одного из ключевых элементов политической власти на его территории; выступающий в рамках капиталистического воспроизводства в качестве одной из форм капитала;
    \item \textit{Капитал}, представляющий собой особое исторически обусловленное общественное отношение, опосредующее эксплуатацию наемного труда в ходе метаморфоза стоимости через специфические организационно-вещные формы капиталистического способа производства и реализующееся в качестве самовозрастающей стоимости в рамках процесса расширенного воспроизводства.
\end{itemize}

Товар заключает внутри себя следующие фундаментальные противоречия:
\begin{itemize}
    \item \textit{Противоречие между овеществленным и абстрактным трудом}, исходящее из существования продукта как такового в качестве формы конкретного труда, не объективированного в подвластной рыночным законам форме товара, являющегося в качестве носителя стоимости выражением абстрактного труда, обусловленного отчуждением от его результатов рабочего как непосредственного производителя;
    \item \textit{Противоречие между потребительной и меновой стоимостью}, заключающееся в его производстве и реализации капиталистом ради приращения стоимости и извлечения прибыли в противоположность приобретению товара покупателем ради его полезных свойств и удовлетворения собственных потребностей;
    \item \textit{Противоречие между стоимостью и ценой}, обусловленное определением стоимости товара на основе объективной величины общественно необходимых для его производства затрат абстрактного труда в противоположность представляющей собой денежную форму стоимости рыночной цене, формирующейся в процессе товарного обмена в рамках рыночного механизма и отклоняющейся от действительной стоимости товара под влиянием различий нормы прибыли и перетока капитала между отраслями, соотношения спроса и предложения, состояния конкурентной среды и прочих факторов ценообразования;
    \item \textit{Противоречие между частным и общественным характером труда}, обусловленное совершаемым вне сознательной координации с потребностями общества индивидуальным трудом множества разрозненных производителей, приобретающим свое общественное выражение исключительно посредством рыночного обмена в рамках сложившейся системы разделения труда;
    \item Противоречие между актами купли и продажи товара, выраженное разложением опосредованного денежными средствами товарного обмена на два разнонаправленных и разорванных во времени и пространстве действия, совершаемых различными субъектами и характеризующихся отсутствием каких-либо гарантий заключения сделки с противоположной стороной.
\end{itemize}

Деньги заключают внутри себя следующие фундаментальные противоречия:
\begin{itemize}
    \item \textit{Противоречие между деньгами как всеобщим эквивалентом и как всеобщей формой стоимости}, обусловленное непрерывным расхождением идеально существующего и независимого от текущего состояния товарного обращения абстрактного выражения стоимости товара в устойчивых номинальных единицах и реально выступающего в непрерывном процессе общественного воспроизводства конкретного ее воплощения в подверженной изменению собственной покупательной способности денежной форме;
    \item \textit{Противоречие между деньгами как средством обращения и как средством платежа}, обусловленное временным разрывом актов отчуждения и оплаты товара, влекущим за собой опосредование всего Воспроизводственного процесса замещающей реальное средство обращения цепочкой взаимно учтенных и менее ликвидных долговых обязательств, требующих в конечном счете своевременного погашения в действительном денежном эквиваленте.
\end{itemize}

Капитал заключает внутри себя следующие фундаментальные противоречия:
\begin{itemize}
    \item \textit{Противоречие между абсолютным накоплением и относительным обесценением капитала}, выраженное стремлением капитала как самовозрастающей стоимости к непрерывному расширению собственной массы и масштаба применения посредством накопления и капитализации прибавочной стоимости, концентрации и централизации средств производства, обуславливающей рост технической оснащенности общественного производства и повышение производительности труда; однако в то же время сокращающей общественно необходимые затраты на производство элементов постоянного капитала и ведущей тем самым к снижению стоимости и частичному или полному обесценению уже функционирующих производственных мощностей.
    \item \textit{Противоречие между постоянным и переменным капиталом}, обусловленное конкуренцией, вынуждающей капиталистов осуществлять постоянную технологическую модернизацию собственных производственных мощностей, увеличивающую производительность труда и снижающую индивидуальные издержки ниже уровня общественно необходимых затрат в целях извлечения добавочной прибыли, однако ведущую по мере распространения новых технологий к снижению доли рабочей силы в качестве переменного капитала и единственного источника прибавочной стоимости, обуславливая тем самым рост органического строения капитала и долгосрочную тенденцию понижения средней нормы прибыли;
\end{itemize}

Капиталистический способ производства заключает внутри себя следующие фундаментальные противоречия:
\begin{itemize}
    \item \textit{Противоречие между пролетариатом и буржуазией}, обусловленное обеспечивающей процесс самовозрастания стоимости общественной природой капитала, заключающей в себе классовый антагонизм частной капиталистической формы эксплуатации человека человеком, выражающийся в противоположности интересов трудящихся и собственников средств производства относительно ее масштабов, определяющих размер заработной платы, интенсивность труда и продолжительность рабочего дня;
    \item \textit{Противоречие между общественным характером производительных сил и частной собственностью на средства производства}, обусловленное отчуждением формирующей подавляющее большинство населения рабочей силы как носителя трудовых навыков и основной производительной силы общества от находящихся в составляющей основу капиталистической формы присвоения результатов труда и эксплуатации человека человеком частной собственности узкого круга непроизводительного класса буржуазии средств производства;
    \item \textit{Противоречие между общественным характером труда и частной формой присвоения его результатов}, обусловленное существованием товара в качестве результата кооперации опирающегося на коллективные достижения в области науки и техники труда множества выходящих за рамки одного предприятия людей, получающих в конечном итоге лишь необходимую для возмещения стоимости собственной рабочей силы часть конечного продукта за вычетом безвозмездно присваеваемой капиталистами созданной трудом рабочих прибавочной стоимости;
    \item \textit{Противоречие между удовлетворением общественных потребностей и извлечением прибыли}, обусловленное организацией производства и распределения преимущественно в рамках частных предприятий на основе интересов максимизации прибыли и накопления капитала, игнорирующих действительные потребности общества в случае невозможности их удовлетворения на чисто капиталистической основе и превращающих необходимые товары и услуги в объекты спекуляции и наживы;
    \item \textit{Противоречие между организацией производства на отдельных фабриках и анархией производства во всем обществе}, находящее свое выражение в стихийной координации капиталистической экономики посредством рыночного механизма в условиях подчинения определяющих пропорции общественного воспроизводства частных управленческих решений интересам извлечения прибыли и вытеснения конкурентов в рамках антагонизма внутриотраслевой и межотраслевой конкуренции различных предприятий, обуславливающей в отсутствие централизованного регулирования процесса воспроизводства на уровне всего общества в целом возникновение регулярных периодических диспропорций и кризисов перепроизводства;
    \item \textit{Противоречие между условиями производства прибавочной стоимости и условиями ее реализации}, возникающее вследствие обусловленного стремлением капиталиста к извлечению дополнительной прибыли роста эксплуатации труда рабочих, сокращающего их относительную платежеспособность и затрудняющего таким образом успешное приращение стоимости в рамках рыночного обмена в ходе метаморфоза капитала между его товарной и денежной формами;
    \item \textit{Противоречие между высокой производственной возможностью капитализма и низкой потребительной способностью народных масс}, выраженное наличием материальной базы в виде развитых производительных сил, технически позволяющей удовлетворить потребности большинства населения в соответствующих товарах и услугах, однако в то же время встречающей на своем пути объективную ограниченность платежеспособного спроса бедностью и пролетарским положением народных масс, обусловленного погоней капитала за максимальной прибавочной стоимостью и нормой прибыли посредством роста производительности и эксплуатацию труда, снижающей долю рабочей силы в стоимости готовой продукции, затрудняющей тем самым ее реализацию и оборачивающейся насильственным разрешением противоречия между относительным перепроизводством и недостаточным на фоне неудовлетворенных потребностей и социальных бедствий общества потреблением в рамках кризисов капиталистической экономики.
\end{itemize}

\subsection{Проявление противоречий капиталистической экономики в рамках кризиса}
В соответствии с основополагающими законами материалистической диалектики, развитие капиталистической экономики опосредовано периодическими кризисами, выступающими не случайными сбоями или внешними нарушениями нормального хода капиталистического воспроизводства, а закономерными этапами диалектического развития капиталистической системы, представляющими собой временное насильственное разрешение ее имманентных противоречий.

Фундаментальные противоречия товара в условиях кризиса капиталистической экономики проявляются следующим образом:
\begin{itemize}
    \item \textit{Противоречие между овеществленным и абстрактным трудом} обостряется и усложняет процесс превращения множества произведенных продуктов труда из формы овеществленного конкретного труда обособленных производителей в признанную обществом форму абстрактного труда товара, однако в то же время составляет основу двойственного характер труда в контексте капиталистического способа производства и в принципе не способно быть снятым в рамках товарной формы труда как таковой;
    \item \textit{Противоречие между потребительной и меновой стоимостью} выражается в утрате произведенными товарами способности быть реализованными в качестве носителей своей исходной меновой стоимости при сохранении их потребительных свойств, частично снимается в рамках опосредованного затяжным накоплением и последующей утилизацией товарных запасов насильственного сокращения избыточного объема товарной массы до потенциально реализуемого уровня;
    \item \textit{Противоречие между стоимостью и ценой} выражается в обусловленном утратой произведенными товарами способности реализовывать заключенную в них меновую стоимость несоответствии текущего уровня товарных цен исходным отношениям стоимости, частично снимается в ходе опосредованного резким обесценением избыточной товарной массы изменения структуры и уровня общественно необходимых затрат и формирования новых отношений стоимости с соответствующим им уровнем товарных цен;
    \item \textit{Противоречие между частным и общественным характером труда} выражается в непризнании рынком овеществленного в товаре частного труда в качестве общественно необходимого, частично снимается в ходе опосредованного банкротствами предприятий и ростом безработицы насильственного исключения избыточной части товаропроизводителей из процесса общественного производства и сокращения тем самым массы признаваемых в качестве общественно необходимых затрат частного труда;
    \item \textit{Противоречие между актами купли и продажи} выражается в разрыве единства опосредованных денежными средствами и обособленных во времени и пространстве актов купли и продажи товара в качестве двух диалектически противоположных составляющих рыночного обмена, частично снимается в рамках опосредованного аннулированием невыполнимых сделок и формированием новой структуры рыночных отношений насильственного сокращения избыточной товарной массы в обращении.
\end{itemize}

Фундаментальные противоречия денег в условиях кризиса капиталистической экономики проявляются следующим образом:
\begin{itemize}
    \item \textit{Противоречие между всеобщим эквивалентом и всеобщей формой стоимости} выражается в обусловленном нарушением метаморфоза капитала между товарной и денежной формами в совокупности с обесценением товаров и капитала остром несоответствии фиксированного номинала денег изменившимся условиям реального выражения стоимости, влекущим за собой рост спроса на деньги и их превращение в абсолютную форму богатства; однако характеризуется сохранением денег как всеобщей формы и как всеобщего эквивалента стоимости, составляет основу двойственного характера денег в контексте капиталистического способа производства и в принципе не способно быть снятым в рамках денежной формы стоимости как таковой;
    \item \textit{Противоречие между средством обращения и средством платежа} выражается в обусловленном разрывом опосредующей товарное обращение цепочки взаимных долговых обязательств и недостатком денег как таковых для их полноценного единовременного покрытия превращении кредитной системы в монетарную в условиях резкого роста спроса на наличные деньги в качестве единственного действительного и окончательного средства расчета, частично снимающееся в ходе опосредованного банкротствами предприятий, принудительным списанием и реструктуризацией долгов насильственного сокращения избыточной массы невыполнимых обязательств и последующего формирования новой обслуживающей товарный оборот структуры долговых связей.
\end{itemize}

Фундаментальные противоречия капитала в условиях кризиса капиталистической экономики проявляются следующим образом:
\begin{itemize}
    \item \textit{Противоречие между абсолютным накоплением и относительным обесценением} обостряется в ходе опосредованного неполной загрузкой и остановкой производственных мощностей в совокупности с банкротствами предприятий абсолютного обесценения избыточной массы перенакопленного капитала, однако характеризующееся сохранением капитала в качестве основополагающей формы производственных отношений и в принципе не способное быть снятым в рамках системы расширенного капиталистического воспроизводства как таковой;
    \item \textit{Противоречие между постоянным и переменным капиталом} выражается в резком насильственном сокращении нормы прибыли вследствие несоответствия ее прежнего уровня текущим пропорциям капиталистического воспроизводства и структуре органического строения общественного капитала с возросшей в ходе ранее осуществленной модернизации производственных мощностей долей постоянного основного капитала, частично снимается в рамках опосредованного удешевлением средств производства вследствие обесценения авансированного капитала и снижением заработной платы в период массовой безработицы временного восстановления нормы прибыли в условиях низкой стоимостной базы в соответствии с повышенным уровнем органического строения капитала обновленных производственных мощностей.
\end{itemize}

Фундаментальные противоречия капиталистического способа в условиях кризиса капиталистической экономики проявляются следующим образом:
\begin{itemize}
    \item \textit{Противоречие между пролетариатом и буржуазией} выражается в резком обострении классового антагонизма вследствие опосредованного массовыми увольнениями и ростом безработицы, падением реальной заработной платы, интенсификацией труда и свертыванием социальных гарантий насильственного переноса издержек обесценения капитала и устранения накопленных диспропорций на трудящиеся массы в ходе стремления капитала к сохранению и восстановлению нормы прибыли посредством усиления эксплуатации и концентрации собственности, однако характеризуется сохранением классовой структуры производственных отношений капитализма и в принципе не способно быть снятым в рамках капиталистического способа производства как такового;
    \item \textit{Противоречие между общественным характером производительных сил и частной собственностью на средства производства} выражается в парализации достигшей высокой степени кооперации, концентрации и взаимозависимости системы общественного производства узкими рамками частнокапиталистического присвоения и рыночной реализации, обуславливающей блокировку общественных производительных сил в форме недозагрузки и остановки мощностей, массового высвобождения рабочей силы и разрушения хозяйственных связей, однако характеризуется сохранением частной собственности на средства производства в качестве основы капиталистической формы эксплуатации и в принципе не способно быть снятым в рамках капиталистического способа производства как такового;
    \item \textit{Противоречие между общественным характером труда и частной формой присвоения его результатов} выражается в усилении эксплуатации трудящихся посредством удешевления рабочей силы в ходе повышения интенсивности труда и падения реальной заработной платы под давлением массовой безработицы в рамках опосредованной массовыми банкротствами предприятий и концентрацией собственности насильственной перестройки условий извлечения и структуры перераспределения сокращающейся вследствие срыва реализации и обесценения капитала массы прибавочной стоимости между различными капиталистами и фракциями капитала, однако характеризуется сохранением частной формы присвоения прибавочного продукта и в принципе не способно быть снятым в рамках капиталистического способа производства как такового;
    \item \textit{Противоречие между удовлетворением общественных потребностей и извлечением прибыли} выражается в опосредованном свертыванием подчиненного логике самовозрастания стоимости выпуска товаров и услуг вследствие массовых банкротств предприятий в совокупности с насильственным освобождением капитала от неприбыльной товарной массы в результате обесценения или физического уничтожения нереализованных товарных запасов и сжатием платежеспособного спроса при сохраняющейся невозможности удовлетворения общественных потребностей на чисто капиталистической основе в период всеобщего срыва реализации сокращении доступа масс к жизненно необходимым благам, однако характеризуется сохранением прибыли как господствующего мотива и критерия производства и в принципе не способно быть снятым в рамках капиталистического способа производства как такового;
    \item \textit{Противоречие между организацией производства на отдельных фабриках и анархией производства во всем обществе} выражается в нарушении планомерности работы предприятий вследствие обострения общих диспропорций капиталистического воспроизводства и срыва рыночной координации, приводящих к разрыву хозяйственных связей, недозагрузке мощностей и хаотической корректировке производственных планов под давлением падения реализации и ликвидности, частично снимается в рамках опосредованного централизацией и концентрацией капитала в совокупности с усилением принудительно-централизующих механизмов капиталистической координации в ходе государственной антикризисной политики насильственного приведения структуры производства и обращения в соответствие с новыми пропорциями капиталистической экономики;
    \item \textit{Противоречие между условиями производства прибавочной стоимости и условиями ее реализации} выражается в разрыве между массой извлекаемой в ходе эксплуатации труда прибавочной стоимости и невозможностью ее превращения в прибыль вследствие нарушения товарно-денежного метаморфоза капитала, частично снимается в ходе опосредованного удешевлением рабочей силы под давлением массовой безработицы и повышением нормы прибавочной стоимости вследствие роста эксплуатации труда временного восстановления ее извлекаемой и реализуемой на сузившейся базе рынка массы;
    \item \textit{Противоречие между высокой производственной возможностью капитализма и низкой потребительной способностью народных масс} выражается в срыве реализации и относительном перепроизводстве товаров на фоне ограниченности потребления и роста неудовлетворенных потребностей общества вследствие пролетарского положения трудящихся, частично снимается в рамках опосредованного обесценением капитала и банкротствами предприятий насильственного устранения нереализованной товарной массы в совокупности с увольнением избыточной части рабочей силы и последующим сокращением объемов выпуска в соответствии с сузившимися рамками платежеспособного спроса.
\end{itemize}